\documentclass[11pt]{article}
\usepackage[margin=27mm]{geometry}
\usepackage{amsmath,amssymb,amsthm,booktabs,longtable,hyperref}
\newtheorem{lemma}{Lemma}
\newtheorem{theorem}[lemma]{Theorem}
\newtheorem{proposition}[lemma]{Proposition}
\newcommand{\PSL}{\operatorname{PSL}}
\newcommand{\PGL}{\operatorname{PGL}}
\newcommand{\F}{\mathbb F}
\newcommand{\N}{\mathcal N}
\title{Local arc transitivity from small stabilizers}
\author{Human Compression Project}
\date{}
\begin{document}\maketitle
\begin{abstract}
We determine the local arc lengths of the configurations in Tables 1--6
of Buekenhout, De Saedeleer and Leemans, \emph{On the rank two geometries
of the groups $\PSL(2,q)$: part II}. The proof uses orders of arc
stabilizers, kernels of small permutation actions, and one argument for
affine subfield covers. Two configurations with stabilizers $S_4$ require
a longer calculation; the same elementary facts about Sylow two-subgroups
settle both. These results concern the specified acting group. They do not
determine the full automorphism group of an incidence graph.
\end{abstract}

\section{Arcs and stabilizers}
Let $X$ be a connected bipartite graph with vertex types $0,1$, and let
$G$ act faithfully, preserve types, and be transitive on edges. Suppose
the degrees $d_0,d_1$ are at least two. An $s$-arc is a sequence
$(x_0,\ldots,x_s)$ of adjacent vertices with $x_{i-1}\ne x_{i+1}$.
The action is locally $s$-arc-transitive if $G_x$ is transitive on the
$s$-arcs starting at $x$, for each vertex $x$. We assume local
two-transitivity, equivalently local $2$-arc-transitivity in this setting.
Write $K_x$ for the kernel of $G_x$ on $\N(x)$.

\begin{lemma}\label{count}
The number of $s$-arcs starting at a vertex of type $i$ is
\[
 a_i(s)=d_i\prod_{j=1}^{s-1}(d_{i+j\bmod2}-1).
\]
Local $s$-arc-transitivity implies $a_i(s)\mid |G_x|$ for both types.
Given local $s$-arc-transitivity, extension to $s+1$ is equivalent to
transitivity of $G_{x_0\cdots x_s}$ on
$\N(x_s)\setminus\{x_{s-1}\}$ for each of the two starting types.
\end{lemma}
\begin{proof}
There are $d_i$ choices for the first edge and, at each later step, all
neighbors except the preceding vertex are allowed. This proves the product
formula. Orbit--stabilizer gives divisibility. To compare two longer arcs,
first carry their initial $s$-arcs to one another; the remaining possible
maps are precisely the stabilizer of that initial arc. This proves the
last assertion, including its converse.
\end{proof}

In particular, the stabilizer of a $2$-arc $(x,y,z)$ is a pointwise
two-point stabilizer in the action of $G_y$ on its neighbors. Its order
is $|G_y|/[d_y(d_y-1)]$. This observation uses the middle vertex;
confusing it with an endpoint would give the wrong order.

\begin{lemma}\label{two}
If $d_1=2$, the graph is locally $3$-arc-transitive.
\end{lemma}
\begin{proof}
For a $2$-arc $(x,y,z)$ with $y$ of type $1$, the subgroup fixing $x$
inside $G_y$ fixes both neighbors of $y$, since there are only two.
Thus $G_{xyz}=G_{yz}$. Local two-transitivity at $z$ says that this
point stabilizer is transitive on $\N(z)\setminus\{y\}$. For the
other starting type there is only one possible extension. Apply
Lemma~\ref{count}.
\end{proof}

For a coset graph with parabolics $M,L$ and intersection $J$, the degrees
are $[M:J],[L:J]$. Connectedness is equivalent to $\langle M,L\rangle=G$.
The action is locally two-transitive exactly when the actions of $M,L$
on their cosets of $J$ are two-transitive. These facts follow directly
by translating an incident pair to the standard pair of cosets. In the
applications, the source's retained constructions supply the subgroups
and these local actions. The argument below does not use the source's
exhaustive classification, which requires a correction elsewhere.

\section{Cases settled by orders and prime cycles}
We use $D_{2m}$ for the dihedral group of order $2m$.
The following table gives all cases with equal parabolic types in which
the obstruction to a third arc is immediate.
\begin{center}
\begin{tabular}{ccrr}
\toprule
$M=L$ & $J$ & $d_0=d_1$ & $a_0(3)$\\\midrule
$A_5$ & $D_{10}$ &6&150\\
$A_5$ & $A_4$ &5&80\\
$A_4$ & $C_3$ &4&36\\
$S_4$ & $S_3$ &4&36\\
$S_3$ & $C_2$ &3&12\\\bottomrule
\end{tabular}
\end{center}
Each last entry exceeds the parabolic order, so the exact local arc
length is two. Reversing the types has no effect on that conclusion.

For degree two at one type, Lemma~\ref{two} proves a third arc. The
following obstructions prove that the exact length is three:
\begin{center}
\begin{tabular}{cccrrr}
\toprule
$M$&$L$&$J$&$d_0$&$a_1(4)=2(d_0-1)^2$&$|L|$\\\midrule
$A_5$&$D_{20}$&$D_{10}$&6&50&20\\
$A_5$&$S_4$&$A_4$&5&32&24\\
$A_4$&$C_6$ or $S_3$&$C_3$&4&18&6\\
$S_4$&$D_{12}$&$S_3$&4&18&12\\
$S_3$&$C_2^2$&$C_2$&3&8&4\\\bottomrule
\end{tabular}
\end{center}

\begin{proposition}\label{prime}
The configurations $(A_5,D_{30},D_{10})$ and $(S_4,D_{18},S_3)$
have exact local arc length three.
\end{proposition}
\begin{proof}
Treat the two cases together, with $p=5$ and $p=3$, respectively.
The degrees are $p+1,3$. The action of $D_{6p}$ on its three cosets
of $D_{2p}$ has kernel the rotation subgroup $C_p$ and quotient $S_3$.
Indeed, writing $D_{6p}=\langle r,t\rangle$ with $|r|=3p$, the
subgroup is $\langle r^3,t\rangle$, and the quotient by
$\langle r^3\rangle$ is the usual degree-three dihedral action.

A $2$-arc ending at the first type therefore has stabilizer $C_p$.
Its local action is nontrivial: the degree-six action of $A_5$ is
faithful by simplicity, and the natural degree-four action of $S_4$
is faithful. On the $p$ available next vertices a nontrivial group
of order $p$ is transitive. A $2$-arc ending at the second type has
stabilizer of order two, since the pointwise two-point stabilizers in
these two actions of $M$ have order two. It cannot lie in the kernel
$C_p$, so it exchanges the two available next vertices. This proves
local $3$-arc-transitivity in both orientations. Finally
$a_0(4)=4p(p+1)>2p(p+1)=|M|$, which excludes a fourth arc.
\end{proof}

\section{Affine covers of a subfield point stabilizer}
Let $K=\F_s$, with $s\ge2$, and let $V$ be a vector space of dimension
two over $K$, with a fixed line $W$. Consider the local configuration
\[
 M=\PGL(2,K),\qquad L=V\rtimes K^*,\qquad
 J=W\rtimes K^*.
\]
The embedding of $J$ in $M$ is its usual point stabilizer on
$\mathbb P^1(K)$. In characteristic two $M=\PSL(2,K)$.
The argument only uses these embedded local actions, so it also applies
when $V$ is a two-dimensional subspace of a larger coefficient field.

\begin{proposition}\label{affine}
Every such coset configuration has exact local arc length three.
\end{proposition}
\begin{proof}
The action of $M$ has degree $s+1$. The action of $L$ on its cosets
of $J$ identifies with the affine action on $V/W$, of degree $s$,
whose kernel is $W$. The quotient is $\operatorname{AGL}(1,K)$,
which is sharply two-transitive.

For a $2$-arc ending at an $M$-vertex, the middle vertex has stabilizer
$L$, so the arc stabilizer is exactly $W$. Inside the endpoint group
$M$, this is the translation subgroup of a point stabilizer. It is
regular on the $s$ other points. In the reverse orientation the arc
stabilizer is a two-point stabilizer of $\PGL(2,K)$, of order $s-1$.
It lies in $J$ and has trivial intersection with the translation group
$W$: a nonidentity translation fixes only the point at infinity.
It therefore projects onto the full multiplicative group of the affine
quotient $L/W$ and is transitive on the $s-1$ available next vertices.
This includes $s=2$, when that last set has one element.

Thus both orientations extend to three arcs. At an $L$-vertex the
number of $4$-arcs is $s^3(s-1)$, whereas $|L|=s^2(s-1)$.
Lemma~\ref{count} excludes local $4$-arc-transitivity.
\end{proof}

This applies to the $A_5,E_{16}:3,A_4$ configurations in Table~1,
to the even-characteristic subfield row of Table~4 where it exists,
and to the $\PGL(2,p^n)$ affine row of Table~6. It also proves the
same local value for the missing two-dimensional scalar covers in the
corrected even-characteristic classification. Their existence and orbit
counts are separate assertions and are not inferred from this local proof.

\section{The two-subgroup calculations in $S_4$}
We record the elementary group facts used below. In its action of degree
three on partitions into two pairs, $S_4$ has kernel the normal four-group
\[
 V=\{1,(12)(34),(13)(24),(14)(23)\}.
\]
The three point stabilizers are its Sylow two-subgroups, each isomorphic
to $D_8$. Distinct such subgroups intersect in $V$. Their three central
involutions are the three distinct nonidentity elements of $V$.
These statements follow by taking the stabilizer of
$\{\{1,2\},\{3,4\}\}$ and conjugating it to the other two partitions.

In $D_8=\langle a,b\mid a^4=b^2=1,\ b a b=a^{-1}\rangle$
there are exactly two four-groups,
$\langle a^2,b\rangle$ and $\langle a^2,ab\rangle$.
They intersect in the center $\langle a^2\rangle$. The unique cyclic
subgroup of order four intersects either in that same center.
If this $D_8$ has index two in $D_{16}$, conjugation by an element
outside it interchanges its two four-groups. To see this, write
$D_{16}=\langle r,b\rangle$, $|r|=8$, and $a=r^2$; conjugation by
$r$ sends one of the two displayed four-groups to the other.

\begin{proposition}\label{d24}
The configuration $(S_4,D_{24},D_8)$ has exact local arc length three.
\end{proposition}
\begin{proof}
Both degrees are three. At an $S_4$-vertex the local kernel is $V$.
At a $D_{24}$-vertex it is $C_4$: in the presentation with rotation
$r$ of order twelve, the subgroup $D_8$ is $\langle r^3,b\rangle$,
and quotienting by $\langle r^3\rangle$ gives $S_3$.
Since both local quotients are sharply two-transitive of degree three,
the stabilizer of a $2$-arc is the local kernel of its middle vertex.
Inside the endpoint edge stabilizer $D_8$, neither $C_4$ nor $V$ is
contained in the other. Its nontrivial local image fixes the incoming
neighbor and interchanges the other two. This proves the third arc in
both orientations.

Choose a $3$-arc $(a,b,c,d)$ with $a,c$ of type $S_4$.
The stabilizer of $(a,b,c)$ is $C_4$ and the stabilizer of the $3$-arc
is its subgroup of order two. Its generator is a square of an element
of order four in $G_a=S_4$, hence is a double transposition in the
normal four-group $K_a$. It fixes every neighbor of $a$. For the reversed
$3$-arc $(d,c,b,a)$ it consequently cannot exchange the two possible
extensions. Local $4$-arc-transitivity fails.
\end{proof}

\begin{proposition}\label{s4pair}
Suppose $G$ is nonabelian simple and is generated by the two parabolics
in a configuration $(S_4,S_4,D_8)$. Its exact local arc length is four.
\end{proposition}
\begin{proof}
Write $V_x=K_x$ at each vertex. If $x,y$ are adjacent and $V_x=V_y$,
this nontrivial four-group is normal in both $G_x$ and $G_y$, and hence
in $G$. This contradicts simplicity. Thus $V_x,V_y$ are the two distinct
four-groups in $J_{xy}=G_x\cap G_y$, and
\[
 V_x\cap V_y=Z(J_{xy})=\langle z_{xy}\rangle.
\]
For a $2$-arc $(a,b,c)$, its stabilizer is $V_b$. Its image on the
neighbors of $c$ has order two, since $V_b\cap V_c$ has order two.
That image fixes $b$ and interchanges the other two neighbors. The
stabilizer of a resulting $3$-arc $(a,b,c,d)$ is
$V_b\cap V_c=\langle z_{bc}\rangle$.

The subgroups $J_{bc}$ and $J_{cd}$ are different Sylow two-subgroups
of $G_c=S_4$, so $z_{bc}\ne z_{cd}$. Both belong to $V_c$, whereas
$V_c\cap V_d=\langle z_{cd}\rangle$. Thus $z_{bc}$ does not belong
to the local kernel $V_d$. It fixes $c$ and exchanges the other two
neighbors of $d$, proving the fourth arc. Finally, $a_i(5)=48>24$,
so a fifth arc is impossible for this specified group.
\end{proof}

\begin{proposition}\label{d16}
The configuration $(S_4,D_{16},D_8)$ has exact local arc length seven.
\end{proposition}
\begin{proof}
Write an alternating path as
\[
 (v_0,e_0,v_1,e_1,v_2,e_2,v_3,e_3),
\]
where the $v_i$ have stabilizer $S_4$ and the $e_i$ have stabilizer
$D_{16}$. A vertex $e_i$ has just two neighbors. An element fixing
$e_i$ and one of them fixes the other; hence the two edge stabilizers
there coincide. Write this common group as $J_i\cong D_8$.
Let $V_i$ be the local four-group at $v_i$.

The index-two subgroup $J_i$ is normal in $G_{e_i}=D_{16}$.
An element of $G_{e_i}\setminus J_i$ exchanges its endpoints and
interchanges the two four-groups in $J_i$, by the calculation above.
Consequently
\[
 V_i\ne V_{i+1},\qquad
 V_i\cap V_{i+1}=Z(J_i)=\langle z_i\rangle.
\]
After the first two edges the stabilizer is $J_0$, of order eight.
At $v_1$ it fixes the incoming neighbor and exchanges the other two,
since its kernel there is $V_1$. The stabilizer after three edges is
$J_0\cap J_1=V_1$, and the fourth edge is forced.

At $v_2$ the image of $V_1$ has order two, since
$V_1\cap V_2=\langle z_1\rangle$. It fixes the incoming neighbor
and exchanges the other two. Thus the fifth step is transitive and
its stabilizer is $\langle z_1\rangle$. The sixth step is forced.
At $v_2$, the two different Sylow subgroups $J_1,J_2$ have different
centers, so $z_1\ne z_2$. Since
$V_2\cap V_3=\langle z_2\rangle$, the element $z_1\in V_2$ is
outside $V_3$. It exchanges the two continuations at $v_3$. This
proves transitivity on all $7$-arcs starting at the $S_4$ type.

This also gives the other starting type. Vertex transitivity on the
$S_4$ half and the proved local transitivity imply that $G$ is transitive
on all $7$-arcs starting in that half. Reversal is a $G$-equivariant
bijection to the $7$-arcs starting in the other half, since seven is odd.
Hence those arcs also form a single $G$-orbit. Two of them with the same
initial vertex can only be mapped to one another by its stabilizer,
which proves the required local transitivity.
At a degree-two vertex, $a_1(8)=32>|D_{16}|=16$, which excludes
an eighth arc.
\end{proof}

\section{Comparison with the fixed source}
The following is the complete assignment of the local arc column. It
does not assert that all accompanying existence or equivalence counts
in the source tables are correct; those are separate source obligations.
\begin{center}\small
\begin{tabular}{clll}
\toprule
Source table&Rows&Exact length&Argument\\\midrule
1&$7,8,10,11,14,15,17,18$&2&equal-type order bound\\
1&$1,6,13$&3&Proposition~\ref{prime}\\
1&$2,3,4$&3&Proposition~\ref{affine}, $s=4$\\
1&$5,9,12,16$&3&degree two and order bound\\
2&$1,2,3,4,5$&3&degree two and order bound\\
2&$6,7$&2&equal-type order bound\\
3&$1,9$&3&degree two and order bound\\
3&$2$&3&Proposition~\ref{prime}\\
3&$3,4$&2&equal-type order bound\\
3&$5$&7&Proposition~\ref{d16}\\
3&$6$&3&Proposition~\ref{d24}\\
3&$7,8$&4&Proposition~\ref{s4pair}\\
4&$1$&3&Proposition~\ref{affine}\\
4&$2$&2&equal-type order bound\\
4&$3,4$&3&degree two and order bound\\
4&$5$&3&Proposition~\ref{prime}\\
5&$1,2$&2&$A_4,A_4,C_3$ order bound\\
6&$1$&3&Proposition~\ref{affine}\\
6&$2$&3&degree two and order bound\\
6&$3$&4&Proposition~\ref{s4pair}\\\bottomrule
\end{tabular}\end{center}

In Section~6 the printed second orientation of the $A_5,D_{30}$
argument repeats types $0,1,0$; it must read $1,0,1$. In the
$S_4,D_{24}$ argument the printed ``$D_{18}$ on cosets of $D_8$''
must read $D_{24}$. Proposition~\ref{d24} retains the obstruction
attributed there to an unnamed referee. The general arc-stabilizer
method is standard and is explicitly attributed by the source to
Leemans's earlier work. The contribution here is the complete uniform
derivation of the table column, especially the two $S_4$ calculations
and their common use of distinct Sylow centers.

The separate assertions on the full automorphism groups of particular
graphs on source p.~386 are not settled by this note. An automorphism
group containing the specified $G$ can have greater arc transitivity;
the $q=9$ example explicitly illustrates that distinction. No upper
bound on such a larger group follows merely because an outer automorphism
of $G$ interchanges two subgroup classes. No full-source completion,
Lean verification, new open-question answer or compression factor is
claimed by this component.

\begin{thebibliography}{9}
\bibitem{source} F.~Buekenhout, J.~De Saedeleer and D.~Leemans,
\emph{On the rank two geometries of the groups $\PSL(2,q)$: part II},
Ars Mathematica Contemporanea \textbf{6} (2013), 365--388,
Tables~1--6 and Section~6.
\bibitem{method} D.~Leemans, \emph{Locally $s$-arc-transitive graphs
related to sporadic simple groups}, Journal of Algebra \textbf{322}
(2009), 882--892; attribution follows the fixed source.
\end{thebibliography}
\end{document}
